% Generated draft: requires mathematical and source review.
\documentclass[11pt]{article}
\usepackage[margin=25mm]{geometry}
\usepackage{fontspec}
\setmainfont{Latin Modern Roman}
\usepackage{amsmath,amssymb,mathtools,unicode-math}
\setmathfont{Latin Modern Math}
\usepackage{enumitem,needspace,xurl,hyperref}
\setlength{\parindent}{0pt}
\setlength{\parskip}{5pt}
\setlength{\emergencystretch}{3em}
\newcommand{\N}{\mathbb{N}}
\newcommand{\Z}{\mathbb{Z}}
\newcommand{\Q}{\mathbb{Q}}
\newcommand{\R}{\mathbb{R}}
\newcommand{\C}{\mathbb{C}}
\newcommand{\F}{\mathbb{F}}
\newcommand{\A}{\mathbb{A}}
\newcommand{\PP}{\mathbb{P}}
\newcommand{\E}{\mathbb{E}}
\newcommand{\eps}{\varepsilon}
\newcommand{\dd}{\,\mathrm{d}}
\begin{document}
\section*{UnsolvedMath ID: 7200082}
\textbf{Source problem code:} catalogue entry
\subsection{Definitions and mathematical statement}
Let $K$ be a tame knot type represented by closed embedded curves in $\R^3$, and let $\operatorname{Cr}(K)$ be its crossing number, the minimum crossing count among generic planar diagrams. In the standard radius-based convention, the ropelength of a $C^{1,1}$ embedding $\gamma$ is
\[
 \operatorname{Rop}(\gamma)=
       \frac{\operatorname{length}(\gamma)}{\operatorname{thickness}(\gamma)},
\]
where thickness is the maximal embedded normal-tube radius, equivalently the curve's reach in the usual convention. Define $L(K)$ as the infimum of this quantity over embeddings of type $K$.

The source asks for the constant in the universal inequality
\[
 L(K)\ge a\,\operatorname{Cr}(K)^{3/4}.
\]
Its research summary identifies the task as finding the optimal constant. A precise extremal expression is
\[
 a_* =\inf_{K:\operatorname{Cr}(K)>0}
       \frac{L(K)}{\operatorname{Cr}(K)^{3/4}}.
\]
The bound with coefficient $a_*$ is required for every knot type, not a coefficient chosen separately for each $K$. The unknot has crossing number zero and contributes no restriction to this ratio. The formula does not assume an extremizing knot realizes the infimum.

The supplied question reports $1.1<a\le10.76$ and the summary says existence of a positive coefficient and the exponent $3/4$ are established, while the sharp coefficient is unknown. The numerical range is retained as supplied evidence, not independently established here. Thickness as radius versus diameter changes the numerical coefficient; the source does not define its normalization, curve class, or minimum-versus-infimum convention. These must be matched before the displayed reported constants can be accepted in the chosen normalization.
\subsection{Short English statement}
Determine the best universal coefficient multiplying crossing number to the power $3/4$ in the minimum-ropelength lower bound. The stated interval $1.1<a\le10.76$ needs the source's precise thickness normalization.

The supplied summary attributes existence of a bound to Diao and improvements to Diao--Ernst and collaborators, but gives no exact paper or theorem. It reports no 2024--2026 determination of the sharp coefficient. This is a supplied partial-progress assessment, subject to the pinned Wikipedia record's primary-source review. The background's unpublished note conflicts with published registry metadata.
\subsection{Sources}
\begin{itemize}
\item Local UnsolvedMath registry, permanent ID 7200082, source code catalogue entry, ropelength/crossing-number constant question. Dataset curation and metadata: Ulam AI and UnsolvedMath Contributors, CC BY 4.0.
\item Wikipedia contributors, \emph{List of unsolved problems in mathematics}, Geometry, bullet 82, pinned revision 1366636767, \url{https://en.wikipedia.org/w/index.php?title=List_of_unsolved_problems_in_mathematics&oldid=1366636767#Geometry}; MediaWiki-source extraction, accessed 2026-07-29.
\item Diao; Diao--Ernst and collaborators, as named in the research summary for the $3/4$ bound and improvements. No complete primary references or precise normalizations are supplied.
\item Source-specific provenance: primary-source status audit NEEDS\_REVIEW; background says unpublished pending audit although published metadata is true. Wikipedia text is CC BY-SA 4.0, with attribution and ShareAlike handling required for release, separately from dataset attribution.
\end{itemize}
\end{document}
